% Section 10.3, from lectures/L10/appendix/b_exercises.md.

\section{Exercises}
\label{c10:sec:exercises}\label{c10:app:b}

\begin{aside}
\textbf{How to work these.} The Reflection exercises are answered on paper, from the chapter. The
Code exercises are programs for the board: build and flash them the way \secref{c2:app:c}
(Microchip Studio) or \secref{c2:app:d} (Linux) describes, or run them in the simulator of
Appendix~\ref{sim:ch} if you have no board yet.

\textbf{How to check your work.} The Reflection exercises are answered in
Appendix~\ref{app:answers}. The Code exercises have no printed solution: the board tells you
whether a program does what the exercise describes. Try each exercise before you read its answer.
\end{aside}

% ------------------------------------------------------------------------------------------
\exerciseset{Concepts}
\label{c10:set:1}

\reflexercise{Only a forward declaration}
\label{c10:ex:1}
Why does \file{gpio.h} only forward-declare \code{gpio_t} (\code{typedef struct gpio gpio_t;})
instead of including the full field list? What would go wrong, from an encapsulation standpoint, if
\code{struct gpio { ... };} were written directly in the header instead?

\reflexercise{When \code{malloc()} fails}
\label{c10:ex:2}
\code{gpio_create()} returns \code{NULL} when \code{malloc()} fails instead of, say, overwriting
the first GPIO it ever created. Why does that matter, and what should calling code do with the
return value?

% ------------------------------------------------------------------------------------------
\exerciseset{Program}
\label{c10:set:2}

\codeexercise{Active-low GPIOs}
\label{c10:ex:3}
Add a second behavior to \secref{c10:a:2}'s driver: \code{static} functions
\code{gpio_write_inverted()}, \code{gpio_read_inverted()}, and \code{gpio_toggle_inverted()} that
treat the pin as active-low (a \code{true} write clears the bit rather than setting it), plus a
second \code{gpio_init_inverted(gpio_t* self, uint8_t pin, gpio_mode_t mode)} that behaves like
\code{gpio_init()} but wires a \code{gpio_t}'s pointers to these instead of the originals. Confirm
the same \code{led1.write(&led1, true)} call site now does the opposite, with no change to any
calling code, only which functions \code{gpio_init_inverted()} wired up.

\codeexercise{An opaque timer}
\label{c10:ex:4}
Make Chapter~\ref{c9:ch}'s \code{timer_t} opaque: move \code{struct timer { ... };} into a
\file{timer.c} file, expose only \code{typedef struct timer timer_t;} plus prototypes for
\code{timer_create()}, \code{timer_clear()}, \code{timer_running()}, \code{timer_start()},
\code{timer_stop()}, \code{timer_toggle()}, and \code{timer_reset()}, from \mbox{\file{timer.h}}.
\code{timer_create(uint32_t timeout_ms, timer_callback_t callback)} should
follow \code{gpio_create()}'s allocation pattern from \secref{c10:a:3}, since an opaque
\code{timer_t} can no longer be declared and passed by address the way Exercise~\ref{c9:ex:6}'s
\code{timer_init(timer_t* self, uint32_t timeout_ms, timer_callback_t callback)} expected.

\codeexercise{A button callback for real}
\label{c10:ex:5}
Wire up \secref{c10:a:4}'s \code{button_t} for real: create one \code{gpio_t} configured as a
button input and another as an LED output, then set \code{on_press} to a function that calls
\code{gpio_toggle()} on the LED. \code{button_check()}'s own code should never mention the LED's
\code{gpio_t} by name. See \secref{c10:a:5} for a fully wired, interrupt-driven version of this
same idea.

\codeexercise{Testing through the stub}
\label{c10:ex:6}
Capstone, harder. \code{gpio_stub_new()} (\secref{c10:a:6}) exists specifically so GPIO-dependent
logic can be exercised without a chip attached. Write a small test that creates a
\code{gpio_interface_t*} via \code{gpio_stub_new()} alone, no \code{gpio_atmega328p_new()}, no
hardware, and confirms \code{write()}/\code{read()}/\code{toggle()} behave correctly purely through
the interface (write \code{true}, read back \code{true}, \code{toggle()} flips it to
\code{false}). Then write a third implementation, one that logs every \code{write()}/%
\code{toggle()} call before updating its own \code{state}, and confirm \file{main.c}'s
edge-detection logic (\secref{c10:a:6}) works against it completely unmodified, just by changing
which \code{gpio_interface_t*} gets created.
